\section{Codimension-two repair is a weighted grid problem}
\label{sec:quotient-grid}

Let $p$ be an odd prime, $N=4p$, and $m=N/2$. Fix a self-dual code
$C\subseteq\mathbb F_p^N$ and a subspace $G\subseteq C$ of dimension
$m-2$. Assume $G$ contains at least five orthogonal sign classes.
Write $b=|S(G)|$ and let $\mathcal P$ be the set of normalized sign
classes whose residues lie in $G^\perp$, with $M=|\mathcal P|$.
The four-dimensional quotient $G^\perp/G$ converts compatibility
search into an explicit weighted grid.

The underlying two-ruling geometry is classical: the split hyperbolic
quadric $Q^+(3,p)$ has $(p+1)^2$ points on two systems of $p+1$
generator lines. For a primary research source stating this structure,
see Edoukou, \emph{The weight distribution of the functional codes
defined by forms of degree 2 on Hermitian surfaces}, Theorem 4.2,
\url{https://www.numdam.org/article/JTNB_2009__21_1_131_0.pdf}.
We give the matrix proof needed here rather than relying on a
classification implementation. No novelty claim is made for the
finite geometry or the resulting reformulation.

\begin{theorem}[Weighted-grid classification]
The self-dual codes containing $G$ correspond to the $2(p+1)$ rows
and columns of a $(p+1)\times(p+1)$ grid. There are nonnegative
integer point weights $w_{ij}$ such that
\[
 M=b+\sum_{ij}w_{ij},\qquad
 |S(C_{\mathrm{row}\ i})|=b+\sum_jw_{ij},\qquad
 |S(C_{\mathrm{col}\ j})|=b+\sum_iw_{ij}.
\]
Every sign class outside $G$ is counted at exactly one grid point,
and belongs to exactly two of these codes. The classes in $G$ belong
to every code.
\end{theorem}
\begin{proof}
The quotient $G^\perp/G$ is nondegenerate of dimension four and has
the two-dimensional totally isotropic subspace $C/G$. Choose a
hyperbolic basis $e_1,e_2,f_1,f_2$, with $e_i\cdot f_j=\delta_{ij}$
and all other basis pairings zero. For explicit construction, choose
dual vectors $g_i$ to the $e_i$ and replace
$g_i$ by $f_i=g_i-\frac12\sum_j(g_i\cdot g_j)e_j$.
The quotient norm of $x_1e_1+x_2e_2+y_1f_1+y_2f_2$ is
$2(x_1y_1+x_2y_2)$. Identify this vector with
\[
 A=\begin{pmatrix}x_1&x_2\\-y_2&y_1\end{pmatrix},
 \qquad \tfrac12\|\bar v\|^2=\det A.
\]
A nonzero isotropic vector is a rank-one matrix $uv^{\mathsf T}$.
Its projective class is determined by a pair
$([u],[v])\in\mathbb P^1(\mathbb F_p)^2$, giving the grid.
The totally isotropic planes are precisely the matrices with fixed
column line $[u]$ or fixed row line $[v]$. Indeed two independent
rank-one matrices span a rank-one plane only when their column lines
or their row lines agree. These are the two rulings, each of size $p+1$.
Their inverse images in $G^\perp$ are exactly the self-dual codes
containing $G$.

Every sign vector has norm $N=0$ in $\mathbb F_p$, so its quotient
is zero or isotropic. The zero quotient contributes $b$ common
classes. At each nonzero projective quotient point count its sign
classes as $w_{ij}$. Membership in the two ruling planes proves
all the displayed formulas.
\end{proof}

\begin{theorem}[Exact compatibility and optimal retained-space repair]
Every pairwise orthogonal subfamily of $\mathcal P$ lies in the
sign set of one of the $2(p+1)$ codes above. Consequently the maximum
size of a real orthogonal sign frame preserving the prescribed
retained rows spanning $G$ is exactly
\[
 \boxed{b+\max\left\{\max_i\sum_jw_{ij},\ \max_j\sum_iw_{ij}\right\}.}
\]
Thus full completion with this retained space is equivalent to a
grid line of weight $N-b$, and finding the optimum requires only
row and column sums after the sign pool has been computed.
\end{theorem}
\begin{proof}
All pool classes have the same sign parity as the five retained
orthogonal rows: an opposite-parity class would have inner-product
magnitude $N/2$ against all five, contradicting Bessel. Thus for
distinct pool classes real orthogonality is equivalent to orthogonality
modulo $p$. In the quotient, two rank-one matrices $uv^{\mathsf T}$
and $st^{\mathsf T}$ are orthogonal precisely when their column
lines or their row lines agree. This follows by polarizing the
determinant: its cross term is, up to the fixed convention,
$\det(u,s)\det(v,t)$.

Classes projecting to the same point are mutually compatible;
common $G$ classes are compatible with everything. A family of grid
points pairwise sharing a row or a column lies in one row or one
column: once two distinct points share a row, a point outside that
row cannot share both of their different columns. The one-point
case is immediate. Thus every compatible family lies in a ruling
plane. Conversely all sign classes in a ruling-plane code form
a real orthogonal frame, because that self-orthogonal code contains
the same five retained rows. Including all common classes proves
the exact maximum and constructively attains it.
\end{proof}

Put $W=M-b$. Double counting the grid gives
\[
 \sum_{D\supseteq G,\ D=D^\perp}|S(D)|=2M+2pb,
 \qquad
 \max_D|S(D)|\ge b+\left\lceil\frac{W}{p+1}\right\rceil.
\]
Also every ruling-line weight is at most $N-b$, hence
$W\le(p+1)(N-b)$. These are exact incidence constraints, not a
guarantee that some line reaches the maximum $N-b$.

There is a similarly explicit row-deletion objective. Let an initial
frame be the full sign set of one line $L_0$, with size $r_0$.
For any other line $L$, put $r_L=b+\sum_{x\in L}w_x$ and
$a_L=|S(C_{L_0})\cap S(C_L)|$. Then $a_L=b$ for a different line
of the same ruling, while $a_L=b+w_x$ for the crossing point $x$
of opposite-ruling lines. The minimum deletion count for an
augmenting exchange preserving $G$ is
\[
 \min_{L:r_L>r_0}(r_0-a_L),
\]
with value infinity if no qualifying line exists. Requiring
$r_L=N$ instead gives the minimum full-completion deletion count.
The compatibility theorem proves the lower bound, and replacing
the old classes outside the intersection by the new line's classes
attains it. Common classes never need to be sacrificed.

\subsection{Two exact quotient grids at order twenty}
The checker embeds the earlier twelve-sign self-dual code over
$\mathbb F_5$ and sorts its normalized sign representatives
lexicographically. The first retained choice uses sorted indices
$0,1,\ldots,7$. It gives $b=8$, $M=48$, and grid weights
\[
\begin{pmatrix}
0&2&1&0&0&1\\1&2&1&2&2&4\\2&4&2&1&1&2\\
1&2&0&0&0&1\\0&1&0&1&0&2\\0&1&0&0&1&2
\end{pmatrix}.
\]
The row sums are $(4,12,12,4,4,4)$ and column sums
$(4,12,4,4,4,12)$. Four of the twelve code extensions therefore
give twenty sign classes, while the other eight give twelve.
The average lower bound gives only fifteen, whereas the exact
weighted-grid optimum is twenty.

The original code is the first row line. The full second column
line shares $8+2=10$ original classes, so its complete twenty-row
frame replaces exactly two old rows. Every one-row deletion of
the original twelve-row frame still spans its ten-dimensional
self-dual code. Consequently such a deletion can acquire no modular
sign class outside the old code, and cannot augment. Thus the
global augmenting and full-completion deletion costs for this
frame are both exactly two, with a finite-field rank proof of
the lower bound and an explicit reconstructed full code for the
upper bound. This conclusion does not depend on a clique search.

For a second choice, retain sorted indices
$(11,10,9,8,7,6,5,4)$, corresponding exactly to the original trap
file's first eight rows. Here $b=9$, $M=51$ and the row-line counts
are $(12,20,12,20,20,12)$; the column-line counts are
$(12,20,12,20,12,20)$. Six of the twelve extensions are full.
The grid has twenty-four occupied points, sum of squared point
weights $96$, and ordered orthogonal distinct-point weight $588$,
agreeing with the separate incidence-moment calculation.

In both cases the twelve reconstructed self-dual bases, their exact
sign masks, the grid labels, and all overlap costs are recorded in
\path{quotients/checks.json}. Every pairwise compatibility relation
in the pool and every output Gram matrix is checked as an integer.
The examples demonstrate exact repair through quotient geometry at
a known order; an all-order theorem forcing enough concentration
on a quotient generator has not been proved.
